\documentclass[11pt]{article}

% Bản tiếng Việt dùng để tác giả đọc và rà soát nội dung. Nhánh T5 chạy bằng
% pdfLaTeX mặc định của Overleaf; nhánh fontspec cho phép kiểm tra bằng XeLaTeX.
\usepackage{iftex}
\ifPDFTeX
  \usepackage[T5]{fontenc}
  \usepackage[utf8]{inputenc}
\else
  \usepackage{fontspec}
  \setmainfont{lmroman10-regular.otf}[
    BoldFont=lmroman10-bold.otf,
    ItalicFont=lmroman10-italic.otf,
    BoldItalicFont=lmroman10-bolditalic.otf
  ]
\fi
\usepackage[a4paper,margin=2.5cm]{geometry}
\usepackage{amsmath}
\usepackage{array}
\usepackage{booktabs}
\usepackage{longtable}
\usepackage{graphicx}
\usepackage{microtype}
\usepackage{xurl}

\renewcommand{\figurename}{Hình}
\renewcommand{\tablename}{Bảng}
\renewcommand{\refname}{Tài liệu tham khảo (References)}
\setlength{\parindent}{1.2em}
\setlength{\parskip}{0.35em}
\raggedbottom
\newcolumntype{L}[1]{>{\raggedright\arraybackslash}p{#1}}

\newcommand{\datasetrepository}{Zenodo}
\newcommand{\datasetdoi}{\url{https://doi.org/10.5281/zenodo.22074838}}
\newcommand{\datasetlicense}{Creative Commons Attribution 4.0 International (CC BY 4.0)}
\newcommand{\coderepository}{\url{https://github.com/MinhTrannnnn/2.5D-}}
\newcommand{\figureplaceholder}[2]{%
  \fbox{\parbox[c][#1][c]{0.92\linewidth}{\centering\textbf{Vị trí đặt hình}\\[0.6em]#2}}%
}

\title{Bộ dữ liệu địa hình 2,5 chiều với ba mức độ khó ghép cặp cho lập kế hoạch đường đi robot mặt đất}
\author{%
  Tran Quang Minh\\
  \small Học viện Công nghệ Bưu chính Viễn thông, Hà Nội, Việt Nam\\
  \small Liên hệ: \texttt{MinhTQ.B22VT351@stu.ptit.edu.vn}
}
\date{}

\begin{document}
\maketitle

\section*{Tóm tắt (Abstract)}

Chúng tôi cung cấp bộ dữ liệu gồm 5.010 bản đồ trường độ cao 2,5 chiều được tạo
bằng thuật toán để nghiên cứu lập kế hoạch đường đi cho robot mặt đất. Dữ liệu được tổ
chức thành 1.670 bộ ba Dễ--Trung bình--Khó ghép cặp thuộc năm họ địa hình. Ba
bản đồ trong mỗi bộ ba dùng cùng một mẫu địa hình gốc và mười cặp điểm đầu--đích. Mỗi bản đồ chứa
lưới độ cao $128\times128$ với kích thước ô 0,25 m, cùng các lớp độ cao mặt đỡ,
độ dốc, độ nhám, chiều cao bậc, chi phí đi qua địa hình và mặt nạ va chạm có xét
kích thước robot. Bản phát hành gồm 50.100 bản ghi bản đồ--nhiệm vụ, 10.020 ảnh
địa hình và ảnh phân tích, 30.060 ảnh đường đi của thuật toán, cùng 90.180 bản
ghi kết quả chạy theo một quy trình cố định cho sáu thuật toán. Cấu trúc tệp
JSON, hạt giống ngẫu nhiên (seed) cố định, mã kiểm tra SHA-256, báo cáo kiểm
định, mã nguồn và môi trường Python với các phiên bản thư viện được cố định hỗ trợ
kiểm tra, tái tạo và tái sử dụng dữ liệu. Tệp NPZ không dùng pickle cung cấp
13 lớp dữ liệu dạng lưới (raster) để tải nhanh hơn; JSON vẫn là tệp dữ liệu gốc.

\section{Bối cảnh \& Tổng quan (Background \& Summary)}

Khi hoạt động ngoài những môi trường trong nhà có cấu trúc, robot mặt đất phải
xem xét nhiều yếu tố hơn một lưới vật cản nằm ngang. Độ nghiêng bề mặt, gián đoạn độ cao
cục bộ và độ nhám có thể quyết định một tư thế hoặc một bước chuyển của robot
có khả thi hay không; đồng thời, thay đổi độ cao có thể làm thay đổi chi phí của
một tuyến đường vốn không va chạm.

Các bộ dữ liệu địa hình tự nhiên hiện có phục vụ nhiều bài toán khác nhau. RUGD,
RELLIS-3D, WildScenes và GOOSE cung cấp quan sát cảm biến thực cùng nhãn ngữ
nghĩa cho môi trường phi cấu trúc\cite{rugd,rellis3d,wildscenes,goose}. ORFD
tập trung vào phát hiện vùng tự do, TartanDrive vào học động lực học xe, còn
STONE vào dự đoán khả năng đi qua địa hình trên lưới voxel 3D
\cite{orfd,tartandrive,stone}. Về đánh giá điều hướng, Verti-Bench đánh giá
khả năng đi qua địa hình rất gồ ghề bằng mô phỏng đa vật lý
\cite{vertibench}, còn MarsPlanBench và MoonPlanBench đánh giá lập kế hoạch toàn
cục trên lưới chiếm chỗ được tạo từ mô hình địa hình số của Sao Hỏa và Mặt Trăng
\cite{planetarypathbench}. BARN gần bài toán điều hướng hơn các bộ dữ liệu nhận
thức, nhưng biểu diễn môi trường tổng hợp bằng lưới vật cản phẳng và không gian
cấu hình thay vì địa hình trường độ cao\cite{barn}. Các tài nguyên này không thể
thay thế lẫn nhau: các bộ dữ liệu nhận thức và động lực học giữ lại độ phức tạp
của cảm biến thực và chuyển động của xe, trong khi BARN cung cấp một môi trường thử nghiệm điều hướng
đầu cuối có kiểm soát. Vì vậy, Bảng~\ref{tab:related_resources} và
Bảng~\ref{tab:related_release} tách riêng cấu trúc dữ liệu với kết quả lập kế
hoạch được phát hành, không xếp hạng chất lượng bộ dữ liệu.

{\footnotesize
  \setlength{\tabcolsep}{3pt}
  \renewcommand{\arraystretch}{1.16}
  \begin{longtable}{@{}L{0.18\linewidth} L{0.16\linewidth} L{0.19\linewidth}
                   L{0.24\linewidth} L{0.15\linewidth}@{}}
  \caption{So sánh cấu trúc với các bộ dữ liệu địa hình tự nhiên và bộ đánh giá điều hướng
  liên quan. Quy mô là đơn vị được phát hành hoặc gán nhãn chính, do đó frame
  cảm biến và bản đồ điều hướng không được coi là các mẫu tương đương. ``k/a''
  nghĩa là không áp dụng; ``k/n'' nghĩa là nguồn được trích dẫn không nêu rõ.}
  \label{tab:related_resources}\\
    \toprule
    Tài nguyên & Nguồn và biểu diễn & Quy mô & Khả năng đi qua địa hình hoặc điều kiện
    robot & Độ khó ghép cặp \\
    \midrule
    \endfirsthead
    \multicolumn{5}{c}{\tablename~\thetable{} (tiếp)}\\
    \toprule
    Tài nguyên & Nguồn và biểu diễn & Quy mô & Khả năng đi qua địa hình hoặc điều kiện
    robot & Độ khó ghép cặp \\
    \midrule
    \endhead
    \midrule
    \multicolumn{5}{r}{Còn tiếp ở trang sau}\\
    \endfoot
    \bottomrule
    \endlastfoot
    RELLIS-3D \cite{rellis3d} & Thực; RGB và LiDAR & 13.556 scan LiDAR; 6.235 ảnh
    (bản đồ/nhiệm vụ: k/a) & Nhãn địa hình ngữ nghĩa; không xét kích thước robot
    & Không \\
    ORFD \cite{orfd} & Thực; RGB và LiDAR & 12.198 cặp đồng bộ (bản đồ/nhiệm vụ:
    k/a) & Ba lớp vùng tự do; không nêu mô hình kích thước robot & Không \\
    RUGD \cite{rugd} & Thực; video RGB & 7.453 frame được gán nhãn (bản đồ/nhiệm
    vụ: k/a) & Nhãn địa hình ngữ nghĩa; không xét kích thước robot & Không \\
    WildScenes \cite{wildscenes} & Thực; ảnh và đám mây điểm LiDAR & 9.306 ảnh
    gán nhãn; 12.148 mẫu đám mây điểm (bản đồ/nhiệm vụ: k/a) & Nhãn ngữ nghĩa;
    không xét kích thước robot & Không \\
    GOOSE \cite{goose} & Thực; ảnh RGB/NIR và đám mây điểm & 10.000 cặp gán nhãn
    trong bản phát hành của bài báo (bản đồ/nhiệm vụ: k/a) & Nhãn ngữ nghĩa;
    không xét kích thước robot & Không \\
    TartanDrive \cite{tartandrive} & Thực; quỹ đạo lái đa phương thức, gồm bản đồ
    độ cao & Khoảng 200.000 tương tác (số bản đồ/nhiệm vụ k/n) & Dữ liệu đáp ứng
    xe, không phải lưới đánh giá khả năng đi qua & Không \\
    STONE \cite{stone} & Thực; camera, LiDAR, radar và voxel 3D & 43 chuỗi; 50.878
    frame (số bản đồ/nhiệm vụ k/n) & Nhãn khả năng đi qua 3D dựa trên hình học
    & Không \\
    BARN \cite{barn} & Tổng hợp; lưới vật cản 2D và không gian cấu hình & 300 bản
    đồ; 300 nhiệm vụ điểm đầu--đích định trước & Không gian cấu hình thích nghi theo
    kích thước robot & Các môi trường độc lập được sắp theo độ khó, không ghép cặp \\
    Verti-Bench \cite{vertibench} & Mô phỏng; bản đồ độ cao 2,5D, ngữ nghĩa và
    vật lý địa hình & 100 bản đồ; 1.000 nhiệm vụ điểm đầu--đích & Chín loại xe và
    tương tác xe--địa hình đa vật lý & Không; các mức độ cao được lấy mẫu độc lập \\
    MarsPlanBench/\newline MoonPlanBench \cite{planetarypathbench} & Được tạo từ dữ liệu sứ mệnh;
    lưới chiếm chỗ 2D từ DTM/DEM hành tinh & Hơn 2.000 bản đồ Sao Hỏa; 36 bản đồ
    Mặt Trăng; số nhiệm vụ k/n & Khả năng đi qua dựa trên ngưỡng dốc; không xét
    kích thước hoặc động lực học của xe tự hành & Không; nhiều ngưỡng dốc không phải
    bộ độ khó ghép cặp \\
    Bản phát hành này & Tổng hợp; trường độ cao 2,5D, tệp JSON và raster NPZ &
    5.010 bản đồ; 16.700 nhiệm vụ ghép cặp duy nhất; 50.100 bản ghi bản đồ--nhiệm
    vụ & Các lớp độ dốc, độ nhám, bậc, chi phí và va chạm có xét kích thước robot & Có:
    1.670 bộ ba Dễ--Trung bình--Khó \\
  \end{longtable}
}

{\footnotesize
  \setlength{\tabcolsep}{4pt}
  \renewcommand{\arraystretch}{1.16}
  \begin{longtable}{L{0.18\linewidth} L{0.44\linewidth} L{0.29\linewidth}}
  \caption{Kết quả lập kế hoạch, cách chia tập cố định và khả năng truy cập của các tài
  nguyên trong Bảng~\ref{tab:related_resources}. ``Giấy phép k/n'' nghĩa là
  nguồn được trích dẫn không nêu giấy phép, không có nghĩa tài nguyên được dùng
  không hạn chế.}
  \label{tab:related_release}\\
    \toprule
    Tài nguyên & Kết quả lập kế hoạch và cách chia tập & Dữ liệu, mã nguồn và giấy phép \\
    \midrule
    \endfirsthead
    \multicolumn{3}{c}{\tablename~\thetable{} (tiếp)}\\
    \toprule
    Tài nguyên & Kết quả lập kế hoạch và cách chia tập & Dữ liệu, mã nguồn và giấy phép \\
    \midrule
    \endhead
    \midrule
    \multicolumn{3}{r}{Còn tiếp ở trang sau}\\
    \endfoot
    \bottomrule
    \endlastfoot
    RELLIS-3D \cite{rellis3d} & Không có cặp đầu--đích hoặc kết quả lập kế hoạch cố định;
    có chia tập & Dữ liệu và mã nguồn; CC BY-NC-SA 3.0 \\
    ORFD \cite{orfd} & Đánh giá phát hiện vùng tự do; không có bộ đường đi; cách chia tập k/n
    & Dữ liệu và mã nguồn; giấy phép k/n \\
    RUGD \cite{rugd} & Không có bộ đường đi; có tập chia cho phân đoạn ảnh & Dữ liệu và
    tập chia; giấy phép k/n \\
    WildScenes \cite{wildscenes} & Không có đường đi; có tập chia cho phân đoạn ảnh;
    mã đánh giá & DOI dữ liệu và mã nguồn; CC BY-NC-SA 4.0 \\
    GOOSE \cite{goose} & Không có bộ đường đi; có tập train/validation/test
    & Dữ liệu: CC BY-SA 4.0; mã nguồn: MIT \\
    TartanDrive \cite{tartandrive} & Không có kết quả lập kế hoạch; có tập train/test
    cho bài toán động lực học & Dữ liệu và mã nguồn; CC BY 4.0 \\
    STONE \cite{stone} & Đánh giá dự đoán trên lưới voxel; không có bộ đường đi cố định;
    cách chia tập: k/n & Bài báo và mã nguồn công khai; dữ liệu dự kiến phát hành;
    CC BY-NC-ND 4.0 \\
    BARN \cite{barn} & Đường đi, thước đo độ khó và kết quả cơ sở; không có
    tập train/test cố định & Dữ liệu và mã nguồn; giấy phép k/n \\
    Verti-Bench \cite{vertibench} & Kết quả mười hệ thống trên toàn bộ nhiệm vụ;
    không dùng nhiệm vụ để huấn luyện & Dữ liệu và mã nguồn; mã nguồn dùng giấy phép MIT, thành phần
    bên thứ ba theo giấy phép riêng \\
    MarsPlanBench/\newline MoonPlanBench \cite{planetarypathbench} & Kết quả thuật toán cổ điển và học
    máy; cách chia tập: k/n & Mã nguồn và dữ liệu công khai; giấy phép k/n \\
    Bản phát hành này & Nhiệm vụ cố định dùng chung, kết quả của sáu thuật toán,
    quy trình cố định và tập train/validation/test theo nhóm ghép cặp & Dữ liệu:
    CC BY 4.0; mã nguồn: MIT \\
  \end{longtable}
}

Trong bối cảnh đó, bản phát hành này bổ sung một phương pháp tạo dữ liệu có
kiểm soát cho bài toán lập kế hoạch robot mặt đất trên lưới: dạng địa hình
quy mô lớn và nhiệm vụ điều hướng được giữ nguyên trong mỗi bộ ba, trong khi độ
khó hình học thay đổi theo quy tắc đã được mô tả. Đơn vị trung tâm là bộ ba
Dễ--Trung bình--Khó dùng cùng mẫu địa hình gốc và cùng các cặp điểm đầu--đích.
Mỗi bản đồ còn đi kèm các lớp phân tích có xét kích thước robot và kết quả đánh
giá của các thuật toán lập kế hoạch với seed cố định. Cấu trúc ghép cặp có kiểm soát là điểm
phân biệt chính so với các môi trường độc lập hoặc chỉ được sắp theo độ khó; nó
cho phép so sánh trên cùng một mẫu gốc trong khi giữ nguyên họ địa hình, hình
dạng địa hình do seed tạo ra và tọa độ nhiệm vụ.

Bộ dữ liệu biểu diễn địa hình bằng trường độ cao 2,5 chiều trên lưới: mỗi vị trí
nằm ngang $(x,y)$ chỉ có một giá trị độ cao $z$ và không có ô thiếu. Cách biểu
diễn này cho phép phân tích hình học trên vùng mà robot chiếm chỗ,
đồng thời vẫn có thể sử dụng trực tiếp cho thuật toán dựa trên lưới hoặc lấy
mẫu. Theo chủ đích,
nó không biểu diễn được nhiều bề mặt tại cùng một tọa độ nằm ngang, chẳng hạn
cầu, phần nhô ra hoặc hang. Năm họ địa hình gồm
\texttt{smooth\_obstacles}, \texttt{rolling}, \texttt{mountain},
\texttt{rugged} và \texttt{plateau}; chúng cung cấp các tổ hợp khác nhau của
khối trơn, địa hình lượn sóng, sống núi, vùng gồ ghề và cấu trúc dạng cao nguyên.

Mỗi mẫu địa hình gốc được sinh bằng một seed riêng rồi chuyển thành ba phiên bản
Dễ, Trung bình và Khó. Ba phiên bản giữ cùng cấu trúc địa hình quy mô lớn và
cùng mười cặp điểm đầu--đích, nhưng có hệ số độ khó và chi tiết bề mặt cục bộ
khác nhau. Vì vậy chúng là dữ liệu ghép cặp, không phải ba bản đồ độc lập. Tiêu
chí chấp nhận độ khó giới hạn tỷ lệ ô đi được, kích thước thành phần liên thông
lớn nhất, tỷ lệ ô đi được thuộc thành phần này, số thành phần liên
thông và mật độ biên va chạm. Quá trình xây dựng còn yêu cầu hệ số cường độ
địa hình $q$ tăng và tỷ lệ ô đi được giảm dần từ Dễ đến Khó.

Bản phát hành có 1.670 bộ ba ghép cặp và 5.010 tệp bản đồ. Mỗi bộ ba đóng góp
mười nhiệm vụ dùng chung, gồm ba nhiệm vụ ngắn, bốn nhiệm vụ trung bình và ba
nhiệm vụ dài. Tổng cộng có 16.700 nhiệm vụ ghép cặp duy nhất và 50.100 bản ghi
bản đồ--nhiệm vụ. Một nhiệm vụ dài trong mỗi bản đồ được đánh dấu là nhiệm vụ
chuẩn và được dùng để đánh giá sáu thuật toán. Các lớp raster, nhiệm vụ, cấu hình,
seed sinh dữ liệu, siêu dữ liệu của từng bản đồ, ảnh minh họa, kết quả thuật toán,
danh sách kiểm tra tính toàn vẹn và quy trình đánh giá có thể đọc bằng máy được phát hành
cùng nhau. Cách tổ chức này phù hợp với mục tiêu chung là làm cho dữ liệu nghiên
cứu dễ tìm và dễ tái sử dụng nhờ siêu dữ liệu rõ ràng và quy trình xử lý được mô
tả\cite{fair}.

Bộ dữ liệu có thể được dùng để so sánh có kiểm soát các phương pháp điều hướng
giữa các mức độ khó địa hình ghép cặp, đánh giá hàm heuristic và chi phí cạnh có
xét địa hình, nghiên cứu độ ổn định của thuật toán ngẫu nhiên dưới ngân sách cố định, hoặc
phát triển chương trình đọc dữ liệu và mô hình học khả năng đi qua địa hình.
Kết quả đánh giá chuẩn được cung cấp để làm mốc tham chiếu và kiểm tra tính khả
thi của đường đi; chúng không nhằm tạo ra một bảng xếp hạng thuật toán có giá trị
phổ quát.

\section{Phương pháp (Methods)}

\subsection{Thiết kế tính toán và quy ước không gian}

Mã nguồn tạo dữ liệu dùng Python 3.9.6, NumPy 2.0.2 và SciPy
1.13.1\cite{numpy,scipy}. Matplotlib 3.9.4 và Pillow 11.3.0 được dùng để tạo
ảnh. Dữ liệu, hình ảnh và kết quả đánh giá trong bản phát hành này được tạo trên
MacBook Pro (mã máy Mac16,1) dùng chip Apple M4 với 10 lõi CPU
(4 lõi hiệu năng và 6 lõi tiết kiệm điện), bộ nhớ hợp nhất 16 GB và macOS
26.5.2. Quá trình sinh dữ liệu dùng tối đa sáu tiến trình; quá trình đánh giá
dùng ba tiến trình. Đây là cấu hình tạo bản phát hành, không phải yêu cầu phần
cứng tối thiểu để đọc dữ liệu. Mỗi bản đồ là lưới vuông $128\times128$ ô với kích
thước ô 0,25 m, tương ứng vùng lấy mẫu danh nghĩa $32\times32$ m. Mảng
raster được lưu theo \texttt{array[y][x]}, trong khi điểm đầu--đích của nhiệm vụ
được lưu theo \texttt{[x,y]}. Một ô lưới được đổi sang tọa độ mét bằng
\begin{equation}
  X=0.25x, \qquad Y=0.25y, \qquad Z=z(y,x),
\end{equation}
trong đó $x$ là chỉ số cột, $y$ là chỉ số hàng và hệ số 0,25 đổi chỉ số ô sang
mét. Hàm $z(y,x)$ trả về độ cao bề mặt tại ô đó, cũng tính bằng mét. Tọa độ tâm ô vì vậy nằm
trong khoảng từ 0 đến 31,75 m trên mỗi trục ngang.

Quá trình sinh dùng seed chính \texttt{7300}. Với chỉ số họ
$f\in\{0,\ldots,4\}$ và chỉ số mẫu trong họ
$i\in\{1,\ldots,334\}$, seed gốc là
\begin{equation}
  s_{f,i}=7300+1{,}000{,}003f+10{,}007(i-1).
\end{equation}
Trong công thức này, số hạng theo $f$ tách seed của năm họ địa hình, còn số
hạng theo $i$ tách các mẫu trong cùng một họ. Nhờ đó, mỗi cặp $(f,i)$ có một
seed gốc riêng và có thể được tạo lại trực tiếp.
Nếu mẫu tạo ra không thỏa các điều kiện chấp nhận, chương trình thử lại tối đa
mười lần theo một thứ tự cố định; lần thử $a$ dùng $s_{f,i}+104{,}729a$. Tệp
bản đồ lưu seed được chọn, seed dùng để tạo nhiệm vụ, cấu hình và giá trị thực
tế của các tham số ngẫu nhiên.

\subsection{Các họ địa hình gốc}

Trong cả năm họ địa hình, mỗi ô đều có đúng một giá trị độ cao; ô bị chặn vẫn
có độ cao chứ không phải là ô thiếu dữ liệu. Trường ngẫu nhiên có tương quan
được tạo bằng cách lọc Gaussian một trường nhiễu có phân phối chuẩn rồi chuẩn
hóa phương sai. Sau đó, chương trình thêm các dạng địa hình hình elip với vị
trí, bán trục, hướng và biên độ lấy từ seed. Cùng seed, cấu hình, phiên bản mã nguồn và các
phiên bản thư viện sẽ tái tạo đúng cùng một bề mặt. Năm họ trong bản phát hành
sử dụng trường ngẫu nhiên ở nhiều thang khác nhau kết hợp với các dạng địa hình
được mô tả bằng công thức.

Gọi $\mathcal{G}_{\sigma}$ là một mảng nhiễu chuẩn độc lập đã được lọc Gaussian
với độ lệch chuẩn $\sigma$ ô, sau đó trừ trung bình và chia cho độ lệch chuẩn
của mảng. Phép lọc dùng điều kiện biên phản xạ. Với \texttt{rolling},
\texttt{mountain} và \texttt{rugged}, nền ngẫu nhiên là
$\beta_f\mathcal{G}_{L/7}+0.42\beta_f\mathcal{G}_{L/18}$, trong đó $L=128$ là
chiều rộng lưới. Tâm đồi được lấy từ $U(-0.36E,0.36E)^2$, với $E=31{,}75$ m là
chiều rộng vật lý được lấy mẫu. Mỗi đồi là một elip Gaussian xoay có bán trục
$r_x\sim U(1.3,4.8)$ m và $r_y\sim U(1.2,4.4)$ m, hướng $U(0,\pi)$ và biên độ
$U(-0.55,1.0)A_{h,f}$. Sống núi dùng bán trục $U(4.0,9.0)$ m và
$U(0.55,1.35)$ m, cùng biên độ $U(0.65,1.0)A_{r,f}$. Khối mesa biên mềm dùng
bán trục $U(1.4,3.5)$ m và $U(1.3,3.1)$ m, biên độ $U(1.3,3.2)$ m và độ mềm
biên logistic $U(0.07,0.16)$. Mỗi mảng gồ ghề dùng bán trục $U(1.5,4.2)$ m và
$U(1.4,3.8)$ m, rồi nhân một trường $\mathcal{G}_{0.8}$ dùng chung với elip
Gaussian và $A_{p,f}$. Cuối cùng, độ nghiêng độc lập $U(-0.018,0.018)$ m/m được
thêm trên hai trục ngang. Tâm của sống núi, mesa và mảng gồ ghề dùng cùng phân
phối như tâm đồi; hướng của chúng tuân theo $U(0,\pi)$.
Bảng~\ref{tab:family-parameters} liệt kê các hằng số riêng theo họ.

\begin{table}[htbp]
  \centering
  \caption{Các hằng số của ba họ địa hình có nhiều loại đặc trưng. Số lượng là cố
  định trong từng họ; vị trí, hướng, bán trục và biên độ tuân theo các phân phối
  nêu trong phần văn bản. Mọi hằng số biên độ đều có đơn vị mét.}
  \label{tab:family-parameters}
  \footnotesize
  \begin{tabular}{lccccc}
    \toprule
    Họ & $\beta_f$ & \shortstack{Đồi\\$(N,A_{h,f})$} &
    \shortstack{Sống núi\\$(N,A_{r,f})$} & \shortstack{Mesa\\$N$} &
    \shortstack{Mảng gồ ghề\\$(N,A_{p,f})$} \\
    \midrule
    \texttt{rolling}  & 0,24 & $(9;2{,}2)$  & $(1;2{,}4)$ & 1 & $(2;0{,}10)$ \\
    \texttt{mountain} & 0,38 & $(12;3{,}8)$ & $(3;4{,}2)$ & 1 & $(2;0{,}14)$ \\
    \texttt{rugged}   & 0,48 & $(13;2{,}8)$ & $(2;3{,}2)$ & 2 & $(6;0{,}26)$ \\
    \bottomrule
  \end{tabular}
\end{table}

Nền của \texttt{smooth\_obstacles} là $0.055\mathcal{G}_{L/12}$. Họ này có
ba khối vòm trong một cụm, đặt tại độ lệch $-3{,}8$, 0 và 3,8 m dọc theo một trục
chọn theo góc $U(0,2\pi)$, với nhiễu vị trí độc lập $U(-1,1)$ m trên mỗi tọa
độ, bán trục $U(3.0,4.7)$ m và
$U(3.0,4.8)$ m, cùng biên độ $U(2.6,3.9)$ m. Hai vòm tách rời dùng cả bán trục
và biên độ từ $U(1.2,1.8)$ m. Tâm cụm nằm trong
$[-0.22R,0.22R]\times[-0.08R,0.28R]$; tâm vòm tách rời nằm trong
$[-0.62R,0.62R]\times[-0.68R,-0.36R]$, với $R=15{,}875$ m là nửa chiều rộng
bản đồ. Hướng vòm tuân theo $U(-0.16,0.16)$ radian.

Nền của \texttt{plateau} là
$0.055\mathcal{G}_{L/12}+0.08\mathcal{G}_{L/16}$. Số plateau được lấy đều từ
\{3,4\}. Tâm plateau phân bố quanh pha $U(0,2\pi)$ với nhiễu góc
$U(-0.32,0.32)$ radian và khoảng cách tâm $U(0.28R,0.58R)$. Hai bán trục tuân
theo $U(2.3,4.8)$ m và $U(2.0,4.3)$ m, hướng theo $U(0,\pi)$, độ cao theo
$U(2.2,3.8)$ m và độ mềm biên logistic theo $U(0.075,0.13)$. Bề mặt gốc của cả
năm họ được lọc Gaussian với $\sigma=0{,}55$ ô rồi dịch giá trị nhỏ nhất về 0.

\subsection{Xây dựng độ khó ghép cặp}

Với mỗi bề mặt gốc, trường chi tiết cục bộ riêng theo họ dùng seed bằng phép XOR
theo bit giữa seed của lần thử và \texttt{0x5F3759DF}. Một mảng nhiễu chuẩn được
lọc Gaussian với độ lệch chuẩn $\sigma_f$ rồi chuẩn hóa. Độc lập với mảng này,
một mảng nhiễu đều $U(0,1)$ được lọc với $\sigma=L/12$; các giá trị dưới phân vị
$\rho_f$ bị loại, phần còn lại được chia cho giá trị lớn nhất rồi lọc thêm một lần
với $\sigma=1{,}2$. Tích của hai trường là $\eta_f$.
Bảng~\ref{tab:detail-parameters} nêu toàn bộ giá trị riêng theo họ.

\begin{table}[htbp]
  \centering
  \caption{Tham số trường chi tiết cục bộ và biên độ phụ thuộc độ khó. $s_f$ là
  hệ số cường độ, $m_f$ là giới hạn biên độ và $\rho_f$ là phân vị dùng để tạo
  mặt nạ trọng số.}
  \label{tab:detail-parameters}
  \small
  \begin{tabular}{lcccc}
    \toprule
    Họ & $\sigma_f$ (ô) & $s_f$ & $m_f$ (m) & $\rho_f$ \\
    \midrule
    \texttt{smooth\_obstacles} & 2,8 & 0,24 & 0,045 & 0,58 \\
    \texttt{rolling}           & 2,3 & 0,38 & 0,070 & 0,52 \\
    \texttt{mountain}          & 1,7 & 0,58 & 0,105 & 0,44 \\
    \texttt{rugged}            & 1,2 & 0,78 & 0,145 & 0,34 \\
    \texttt{plateau}           & 2,8 & 0,28 & 0,055 & 0,58 \\
    \bottomrule
  \end{tabular}
\end{table}

Bề mặt ứng viên cho độ khó được tạo theo
\begin{equation}
  z_d(x,y)=q z_0(x,y)+\alpha_f(q)\,\eta_f(x,y),
\end{equation}
trong đó $z_0$ là bề mặt gốc. Hệ số $q$ làm tăng hoặc giảm các dạng địa hình
quy mô lớn; thành phần $\alpha_f(q)\eta_f$ bổ sung chi tiết gồ ghề cục bộ.
Biên độ của thành phần chi tiết được tính bằng
\begin{equation}
  \alpha_f(q)=\min\!\left\{s_f\left[0.010q+
  0.060\max(0,q-0.75)^{1.24}\right],\ m_f\right\}.
\end{equation}
Hàm $\min$ giới hạn biên độ ở $m_f$, vì vậy độ gồ ghề cục bộ tăng theo $q$
nhưng không vượt quá mức tối đa của từng họ địa hình.
Trong mỗi khoảng độ khó $[q_{\min},q_{\max}]$, chương trình đánh giá đúng 34
giá trị cố định theo công thức
\begin{equation}
  q_k=q_{\min}\left(\frac{q_{\max}}{q_{\min}}\right)^{k/33},
  \qquad k=0,\ldots,33,
\end{equation}
đúng như hàm \texttt{numpy.geomspace}. Nói cách khác, các giá trị được đặt cách
đều trên thang logarit, thay vì cách đều trên thang tuyến tính. Các khoảng chung là
0,18--1,05, 0,48--2,40 và 0,90--4,80 cho Dễ, Trung bình và
Khó. Plateau dùng 0,18--1,25, 1,20--3,20 và 2,80--6,00 vì mặt trên phẳng của
khối cao nguyên vẫn có thể là tư thế robot hợp lệ ngay cả khi cạnh bao quanh không thể đi
qua. Nếu $q<0.62$, ứng viên còn được lọc với $\sigma=1.15(0.62-q)$ ô. Giá trị
nhỏ nhất được dịch về 0 và độ cao được làm tròn sáu chữ số thập phân trước khi
tính các lớp được suy ra và các chỉ số chấp nhận. Trường
\texttt{relief\_severity} chính là hệ số $q$ đã chọn, không phải khoảng
max--min độ cao thực tế.

Việc chọn ứng viên dựa trên lưới va chạm đã xét robot, được mô tả bên dưới.
Bảng~\ref{tab:difficulty-acceptance} trình bày toàn bộ miền chấp nhận và ngưỡng
mật độ biên. Trong mỗi bộ ba được chấp nhận, hệ số $q$ tăng nghiêm ngặt và tỷ
lệ ô đi được giảm hơn 0,01 sau mỗi lần tăng độ khó. Các tổ hợp ứng viên được xếp
theo độ lệch so với tỷ lệ ô đi được mục tiêu; chương trình xét tối đa 12 tổ hợp đầu và giữ tổ hợp có
điểm thấp nhất trong số các tổ hợp tạo được toàn bộ nhiệm vụ dùng chung.

\begin{table}[htbp]
  \centering
  \caption{Tiêu chí chấp nhận độ khó dùng khi sinh địa hình. Ba dòng đầu áp
  dụng cho bốn họ địa hình; plateau dùng miền đã hiệu chỉnh riêng. Phần trên
  ghi các dải hiệu chỉnh, phần dưới ghi các ràng buộc cấu trúc tương ứng.}
  \label{tab:difficulty-acceptance}
  \small
  \setlength{\tabcolsep}{5pt}
  \begin{tabular}{llccc}
    \toprule
    Phạm vi & Độ khó & $q_{\min}$--$q_{\max}$ &
    \shortstack{Tỷ lệ ô\\đi được} & Mục tiêu \\
    \midrule
    Chung   & Dễ          & 0,18--1,05 & 0,72--0,94 & 0,84 \\
    Chung   & Trung bình  & 0,48--2,40 & 0,54--0,76 & 0,65 \\
    Chung   & Khó         & 0,90--4,80 & 0,40--0,65 & 0,52 \\
    Plateau & Dễ          & 0,18--1,25 & 0,80--0,90 & 0,84 \\
    Plateau & Trung bình  & 1,20--3,20 & 0,73--0,82 & 0,78 \\
    Plateau & Khó         & 2,80--6,00 & 0,69--0,76 & 0,73 \\
    \bottomrule
  \end{tabular}

  \medskip
  \begin{tabular}{llcccc}
    \toprule
    Phạm vi & Độ khó & $\lambda_{\min}$ &
    \shortstack{Số thành phần\\tối đa} & $\psi_{\min}$ &
    $\delta_{B,\max}$ \\
    \midrule
    Chung   & Dễ          & $\geq0{,}65$ & 40 & 0,90 & 0,080 \\
    Chung   & Trung bình  & $\geq0{,}45$ & 60 & 0,84 & 0,085 \\
    Chung   & Khó         & $\geq0{,}34$ & 75 & 0,80 & 0,090 \\
    Plateau & Dễ          & $\geq0{,}70$ & 25 & 0,90 & 0,060 \\
    Plateau & Trung bình  & $\geq0{,}62$ & 30 & 0,90 & 0,060 \\
    Plateau & Khó         & $\geq0{,}55$ & 35 & 0,88 & 0,060 \\
    \bottomrule
  \end{tabular}
\end{table}

Với lưới va chạm $H\times W$ ký hiệu $G$, đặt
$F=\{(x,y):G(y,x)=0\}$ và gọi $L_F$ là thành phần liên thông lớn nhất theo đúng
đồ thị chuyển động được định nghĩa bên dưới. Ba chỉ số được lưu là
\begin{equation}
  \phi=\frac{|F|}{HW},\qquad
  \lambda=\frac{|L_F|}{HW},\qquad
  \psi=\frac{|L_F|}{|F|}=\frac{\lambda}{\phi},
\end{equation}
trong đó $\phi$ là tỷ lệ ô đi được trên toàn bản đồ, $\lambda$ là tỷ lệ ô thuộc
thành phần liên thông lớn nhất trên toàn bản đồ, còn $\psi$ là tỷ lệ ô đi được
thuộc thành phần đó. Vì vậy, $\phi$ đo tổng diện tích có thể đi qua, còn
$\lambda$ và $\psi$ cho biết vùng đi được có tập trung thành một vùng lớn hay
bị chia nhỏ.

Mật độ biên va chạm là tỷ lệ các cặp ô kề nhau có trạng thái khác nhau: một ô
đi được và một ô bị chặn. Chỉ các cặp ô kề ngang và dọc được tính:
\begin{equation}
  \delta_B=\frac{
  \sum_{y=0}^{H-1}\sum_{x=0}^{W-2}[G_{y,x}\ne G_{y,x+1}]+
  \sum_{y=0}^{H-2}\sum_{x=0}^{W-1}[G_{y,x}\ne G_{y+1,x}]}
  {H(W-1)+(H-1)W}.
\end{equation}
Dấu ngoặc vuông là hàm chỉ thị: nhận giá trị 1 khi điều kiện bên trong đúng và
0 khi điều kiện sai. Giá trị $\delta_B$ cao hơn biểu thị nhiều đường biên giữa
vùng đi được và vùng bị chặn hơn. Phép tính nhanh thành phần nhị phân tám hướng
cho $\widetilde\lambda$ để lọc nhanh ban đầu; các ứng viên còn lại sau đó được
tính lại bằng đồ thị chuyển động đầy đủ. Với khoảng tỷ lệ ô đi được $[\ell,u]$,
mục tiêu $t$ và tỷ lệ thành phần tối thiểu $\lambda_{\min}$, điểm ứng viên chính
xác là
\begin{equation}
  S=|\phi-t|+4\{[\ell-\phi]_+ +[\phi-u]_+\}
  +3[\lambda_{\min}-\widetilde\lambda]_+,
\end{equation}
trong đó $[x]_+=\max(0,x)$. Hạng tử thứ nhất đo độ lệch so với tỷ lệ ô đi được
mục tiêu. Hạng tử thứ hai phạt ứng viên nằm ngoài khoảng $[\ell,u]$, còn hạng
tử thứ ba phạt ứng viên có thành phần liên thông lớn nhất quá nhỏ. Trọng số 4
và 3 làm cho hai vi phạm này ảnh hưởng mạnh hơn sai lệch thông thường so với
mục tiêu. Ứng viên được xếp theo $S$, lọc bằng mọi ngưỡng
trong Bảng~\ref{tab:difficulty-acceptance}, rồi giữ sáu ứng viên hợp lệ đầu tiên
cho mỗi độ khó. Trước phép tính đồ thị đầy đủ, bước lọc nhanh loại ứng viên
có $\widetilde\lambda<0.95\lambda_{\min}$; ngưỡng $\lambda_{\min}$ chính xác
vẫn được áp dụng sau khi tính lại thành phần liên thông. Điểm của một bộ ba là
$\sum_d S_d$.

\subsection{Phân tích địa hình có xét robot}

Robot tham chiếu có vùng chiếm chỗ (footprint) hình tròn với bán kính 0,45 m
trên mặt phẳng ngang. Với kích thước ô 0,25 m, bán kính trên lưới là
\begin{equation}
  R_c=\left\lceil\frac{0.45}{0.25}\right\rceil=2,
  \qquad
  \mathcal{D}=\{(i,j)\in\mathbb{Z}^2:i^2+j^2\leq R_c^2\}.
\end{equation}
Mặt nạ vì vậy có đúng 13 ô trong cửa sổ $5\times5$: ô tâm, bốn ô cách tâm một
ô theo phương ngang hoặc dọc, bốn ô chéo $(\pm1,\pm1)$ và bốn ô cách tâm hai ô
theo phương ngang hoặc dọc. Cách làm tròn lên này tạo thêm biên an toàn: các điểm ngoài cùng trên
trục cách tâm 0,50 m dù bán kính danh nghĩa là 0,45 m. Cấu hình không được suy
ra từ một robot vật lý cụ thể. Bốn ngưỡng gồm bán kính 0,45 m, độ dốc 28 độ, độ nhám 0,10 m và bậc
cao 0,65 m được cố định trước khi sinh bản phát hành để định nghĩa một robot
tham chiếu và tạo được cả vùng đi được lẫn vùng bị chặn trên cả năm họ địa
hình; chúng không phải giới hạn cơ học phổ quát của robot mặt đất. Đường kính
danh nghĩa là 0,90 m. Lớp độ cao thô được phát hành để người dùng tính lại mọi
lớp có xét vùng chiếm chỗ
với hình học và ngưỡng phù hợp cho robot khác. Một mặt phẳng bình phương tối
thiểu được khớp trên cửa sổ tròn tại
mỗi vị trí tâm robot ứng viên. Giá trị của mặt phẳng tại tâm được lưu dưới tên
\texttt{support\_elevation}; độ lớn của gradient ngang được đổi sang độ và lưu
trong trường
\texttt{footprint\_slope\_degrees}. Độ nhám là căn bậc hai của trung bình bình phương
phần dư độ cao so với mặt phẳng. Chiều cao bậc cục bộ là giá trị lớn nhất của
chênh lệch độ cao tám hướng tại các ô trung tâm nằm trong mặt nạ số hóa. Một lớp độ dốc từ
gradient điểm sau khi làm trơn nhẹ cũng được lưu để phân tích mô tả, nhưng
không được dùng để xác định lưới ô đi được hoặc chi phí đường đi.

Cụ thể, gọi $\mathcal{W}=\{(u_i,v_i,z_i)\}_{i=1}^{N}$ là các mẫu trong cửa sổ
tròn, trong đó $(u_i,v_i)$ là độ lệch theo mét so với tâm cửa sổ.
Do mặt nạ tròn đối xứng quanh tâm, mặt phẳng bình phương tối thiểu
$\widehat z_i=\bar z+a u_i+b v_i$ được tính bởi
\begin{equation}
  \bar z=\frac{1}{N}\sum_{i=1}^{N}z_i,\qquad
  a=\frac{\sum_i u_i z_i}{\sum_i u_i^2},\qquad
  b=\frac{\sum_i v_i z_i}{\sum_i v_i^2}.
\end{equation}
Ở đây, $\bar z$ là độ cao trung bình và cũng là giá trị của mặt phẳng tại tâm;
$a$ và $b$ là độ nghiêng của mặt phẳng theo hai phương ngang. Từ hai hệ số độ
nghiêng này, độ dốc trên vùng chiếm chỗ và độ nhám RMS
lần lượt là
\begin{equation}
  \theta=\frac{180}{\pi}\arctan\!\left(\sqrt{a^2+b^2}\right),\qquad
  r=\sqrt{\frac{1}{N}\sum_{i=1}^{N}
  \left(z_i-\bar z-a u_i-b v_i\right)^2}.
\end{equation}
Trong đó, $\theta$ chỉ đo độ nghiêng chung của mặt phẳng đỡ, còn $r$ đo phần
dao động độ cao còn lại sau khi đã loại độ nghiêng này.
Trong cài đặt, các tổng cục bộ của $z$, $z^2$, $uz$ và $vz$ được đánh giá đồng
thời tại mọi tâm ô bằng tích chập hai chiều của
\texttt{scipy.ndimage}; gradient tại từng ô, lọc Gaussian và giá trị lớn nhất cục bộ cũng được
tính trên toàn mảng NumPy. Cách vector hóa này tránh vòng lặp Python qua từng
ô khi phân tích 82.083.840 vị trí của 5.010 bản đồ. Chiều cao bậc dùng tám mảng
láng giềng đã dịch và một bộ lọc lấy giá trị lớn nhất trên vùng chiếm chỗ.

Với độ dốc $\theta$, độ nhám $r$ và chiều cao bậc $h$, chi phí đi qua địa hình
không có đơn vị được tính bằng
\begin{equation}
  C_{\mathrm{trav}}=
  0.45\frac{\theta}{28^\circ}+
  0.25\frac{r}{0.10\ \mathrm{m}}+
  0.30\frac{h}{0.65\ \mathrm{m}}.
\end{equation}
Ba tỷ số lần lượt so sánh độ dốc, độ nhám và chiều cao bậc với ngưỡng tương
ứng. Các trọng số 0,45, 0,25 và 0,30 quy định mức đóng góp của ba yếu tố vào
chi phí tổng.
Đặt $B(x,y)$ là biến logic cho một tư thế tâm robot. Quy tắc phân loại chính
xác là
\begin{equation}
  B=(\theta>28^\circ)\ \lor\ (r>0.10\ \mathrm{m})\ \lor\
  (h>0.65\ \mathrm{m})\ \lor\ (C_{\mathrm{trav}}\geq1).
\end{equation}
Do dùng phép ``hoặc'', tư thế bị chặn nếu vi phạm bất kỳ ngưỡng riêng lẻ nào
hoặc nếu chi phí tổng đạt 1. Điều này ngăn một tư thế có nhiều yếu tố gần ngưỡng
được coi là đi được chỉ vì từng yếu tố riêng vẫn chưa vượt ngưỡng.
Tư thế đi được khi và chỉ khi $B$ sai và tâm không thuộc dải biên. Cụ thể, tâm
bị chặn khi $x<R_c$, $x\geq W-R_c$, $y<R_c$ hoặc $y\geq H-R_c$. Mặt phẳng
được khớp bằng phép tích chập với dữ liệu phản xạ ở biên mảng; khi
tính chiều cao bậc, các mảng dịch và phép lấy giá trị lớn nhất dùng giá trị biên
gần nhất để mở rộng mảng. Cuối cùng, chương trình đánh dấu trực tiếp dải hai ô
ở mép bản đồ là vùng bị chặn.
Không cần làm giãn vật cản thêm để tái tạo cấu hình tham chiếu; nếu thêm biên an
toàn hoặc đổi hình học robot, người dùng đang tạo một cấu hình điều hướng khác.

\subsection{Mô hình chuyển động và hàm mục tiêu đường đi}

Kiểm tra liên thông và tuyến tham chiếu dùng chuyển động tám hướng giữa tâm ô,
không cho cắt góc chéo. Một bước chuyển bị loại nếu một trong hai đầu bị chặn,
đường chéo đi giữa các láng giềng vuông góc bị chặn, hoặc tỷ số chênh cao tuyệt
đối trên quãng đường ngang, tính từ \texttt{support\_elevation}, vượt
$\tan(28^\circ)$. Với chiều dài ngang $d$, thay đổi độ cao mặt đỡ $\Delta z$ và
trọng số độ dốc $w=3$, chi phí
cạnh chung là
\begin{equation}
  c=\sqrt{d^2+\Delta z^2}
  \left[1+w\left(\frac{|\Delta z|}{d}\right)^2\right].
\end{equation}
Thừa số căn bậc hai là chiều dài ba chiều của bước chuyển. Thừa số trong ngoặc
làm chi phí tăng theo bình phương độ dốc; vì vậy hai bước có cùng chiều dài ba
chiều vẫn có thể có chi phí khác nhau nếu độ dốc khác nhau.
Hàm mục tiêu toàn tuyến là tổng $c$ trên mọi bước chuyển. Quy tắc hợp lệ của
chuyển động được dùng chung khi kiểm tra liên thông, chạy sáu thuật toán và kiểm
tra đường đi. Chi phí $c$ được dùng để xếp hạng cặp điểm đầu--đích, tối ưu Dijkstra
và A*, gán trọng số cạnh PRM, chọn và nối lại nhánh trong RRT*, đồng thời báo cáo
chi phí cuối cùng cho mọi đường thành công. BFS chỉ tối thiểu hóa số bước;
RRT-Connect mở rộng cây theo khoảng cách Euclid và các kiểm tra va chạm/độ dốc,
nên $c$ chỉ được tính sau đó để báo cáo đường tìm được.

\subsection{Xây dựng nhiệm vụ điểm đầu--đích (Start--Goal)}

Nhiệm vụ được chọn trên bản đồ Khó của mỗi bộ ba rồi sao chép sang bản đồ
Dễ và Trung bình. Điểm đầu và điểm đích ứng viên nằm trong thành phần liên thông
hợp lệ lớn nhất; thành phần này được tìm bằng phép duyệt BFS dùng hàng đợi trên
đồ thị tám hướng với đúng quy tắc chuyển động đã công bố. Các điểm ứng viên
thông thường cách biên ít nhất 1,5 m, tương đương sáu ô. Tối đa 30 ô gốc được
chọn: dùng toàn bộ khi có không quá 30 ứng viên; nếu có nhiều hơn, chương trình
lấy ngẫu nhiên đúng 30 ô không lặp bằng seed đã định. Thuật toán Dijkstra\cite{dijkstra} được chạy từ từng ô gốc
để tính khoảng cách có trọng số đến các ứng viên khác; các khoảng cách này được
dùng để xếp hạng các cặp điểm đầu--đích theo hàm mục tiêu chung.

Theo quy trình chọn mười nhiệm vụ, ba cặp được chọn từ khoảng phân vị khoảng cách
0,20--0,42, bốn cặp từ 0,45--0,72 và ba cặp từ 0,75--0,98; chúng lần lượt nhận
nhãn \texttt{short}, \texttt{medium} và \texttt{long}. Nhiệm vụ dài nhất được
đánh dấu là nhiệm vụ chuẩn để minh họa và đánh giá. Một bộ ba chỉ được chấp nhận
nếu mọi cặp đã chọn vẫn liên thông theo mô hình chuyển động công bố ở cả ba mức
độ khó.

Nói ngắn gọn, chương trình xếp các cặp từ gần đến xa, lấy các đoạn tương ứng
với ba khoảng phân vị rồi chọn các cặp cách đều trong từng đoạn. Các quy tắc
dưới đây chỉ làm rõ cách xử lý chính xác khi có giá trị bằng nhau, cặp trùng
hoặc không đủ cặp trong một đoạn.

Để xử lý chính xác trường hợp trùng và thiếu, cặp điểm đầu--đích được lưu dưới dạng
cặp không thứ tự rồi xếp trước hết theo khoảng cách có trọng số, sau đó theo thứ
tự từ điển của hàng và cột. Với danh sách có độ dài $n$ và khoảng phân vị
$[p_l,p_u]$, gọi đoạn danh sách tương ứng là $\mathcal{B}$. Đoạn này bắt đầu tại
$\min(n-1,\lfloor p_l(n-1)\rfloor)$ và kết thúc
ngay trước $\min(n,\max(i_l+1,\lfloor p_u n\rfloor))$, trong đó $i_l$ là chỉ
số bắt đầu. Với một khoảng cần chọn $m$ cặp, thuật toán lấy $m$ vị trí cách đều
từ 0 đến $|\mathcal{B}|-1$, làm tròn theo quy tắc số chẵn gần nhất của
Python/NumPy rồi duyệt vòng về phía trước đến cặp chưa dùng đầu tiên. Nếu đoạn
phân vị có ít hơn $m$ ứng viên, toàn bộ danh sách đã xếp hạng sẽ thay thế đoạn
đó. Nếu ba khoảng vẫn cung cấp ít hơn
mười cặp duy nhất, các cặp chưa dùng được lấy từ toàn bộ danh sách theo khoảng
cách giảm dần và gán nhãn \texttt{long}. Lần sinh hiện tại thất bại nếu vẫn
không tạo được mười cặp khác nhau. Nếu có ít hơn hai ô thỏa điều kiện cách biên,
toàn bộ thành phần lớn nhất được dùng; nếu nó cũng có ít hơn hai ô,
lần sinh thất bại.

\subsection{Thuật toán tạo bộ ba hoàn chỉnh}

\noindent\textbf{Thuật toán 1: Sinh một nhóm địa hình ghép cặp.}
Với họ $f$ và chỉ số trong họ $i$:
\begin{enumerate}
  \item Tính $s_{f,i}$ và với lần thử $a=0,\ldots,9$, đặt seed của lần thử là
  $s_{f,i}+104{,}729a$. Sinh một bề mặt gốc $z_0$ và một trường chi tiết
  $\eta_f$ từ seed này.
  \item Với mỗi độ khó, đánh giá đủ 34 giá trị $q_k$. Với từng giá trị, tạo
  $z_d$, thực hiện lọc có điều kiện và làm tròn sáu chữ số, tính các lớp và chỉ số
  có xét vùng chiếm chỗ, tính $S$, rồi giữ sáu ứng viên hợp lệ đầu tiên sau khi xếp
  theo $S$.
  \item Liệt kê mọi tổ hợp có thể tạo từ ba danh sách ứng viên. Loại tổ hợp nếu không
  thỏa $q_{\mathrm{easy}}<q_{\mathrm{medium}}<q_{\mathrm{hard}}$ và
  $\phi_d-\phi_{d+1}>0.01$. Xếp các tổ hợp còn lại theo $\sum_d S_d$.
  \item Với tối đa 12 tổ hợp đầu, tạo mười nhiệm vụ trên bản đồ Khó bằng
  seed thu được từ phép XOR theo bit giữa $(s_{f,i}+104{,}729a)$ và
  \texttt{0x6A09E667}, theo các quy tắc xác định ở trên. Loại tổ hợp nếu
  điểm đầu hoặc điểm đích của bất kỳ nhiệm vụ nào mất liên
  thông trong một trong ba bản đồ.
  \item Trả về tổ hợp hợp lệ đầu tiên và lưu ba bản đồ với cùng các
  nhiệm vụ. Nếu không có tổ hợp hợp lệ, chuyển sang lần thử $a+1$; báo lỗi sinh
  dữ liệu sau mười lần thử không thành công.
\end{enumerate}

\subsection{Đánh giá các thuật toán tham chiếu}

Bản phát hành cung cấp kết quả đánh giá theo một quy trình cố định cho tìm kiếm theo chiều
rộng (BFS),
Dijkstra\cite{dijkstra}, A*\cite{astar}, probabilistic roadmaps (PRM)
\cite{prm}, RRT-Connect\cite{rrtconnect} và RRT*\cite{rrtstar}. Phép đánh giá dùng
nhiệm vụ dài chuẩn trên toàn bộ 5.010 bản đồ. Phạm vi này được cố định trước
khi chạy để có một nhiệm vụ dài được định nghĩa thống nhất cho mỗi bản đồ và
bao phủ đầy đủ ba mức độ khó ghép
cặp theo cùng một quy trình tính toán. Đây là phép đánh giá tham chiếu để kiểm định,
không phải phép đánh giá toàn diện cả mười nhiệm vụ. Toàn bộ mười nhiệm vụ vẫn
được phát hành và chương trình dòng lệnh hỗ trợ chạy tất cả nhiệm vụ mà không
cần sinh lại địa hình; kết quả trong bài không được khái quát cho chín nhiệm vụ
còn lại. BFS là mốc tham chiếu không trọng số, tối thiểu hóa số bước chuyển;
Dijkstra và A* tối ưu hàm mục tiêu có trọng số chung. Các thuật toán không ngẫu nhiên
chạy một lần trên mỗi bản đồ. PRM, RRT-Connect và RRT* chạy với năm seed cố định;
hệ thống không tự động chạy thêm ngoài năm seed đã công bố.

PRM lấy mẫu 500 ô tự do và xét 12 láng giềng gần nhất cho mỗi nút. RRT-Connect
có ngân sách tối đa 5.000 vòng lặp, bước mở rộng bốn ô và dừng khi hai cây kết
nối hoặc khi hết ngân sách. RRT* chạy đủ 6.000 vòng lặp, dùng bước mở rộng ba ô,
bán kính nối lại bảy ô và xác suất lấy mẫu Goal 0,08, sau đó trả về kết nối Goal
hợp lệ có chi phí nhỏ nhất đã tìm thấy. Mỗi cạnh của các thuật toán lấy mẫu được
chuyển thành chuỗi ô bằng thuật toán Bresenham tám hướng đối xứng và chỉ được
chấp nhận khi mọi bước trong chuỗi vượt
qua kiểm tra va chạm, góc và độ dốc chung.

Mỗi lần chạy lưu trạng thái thành công, tính hợp lệ của tuyến, thời gian chạy, chiều dài
phẳng và ba chiều, chi phí có trọng số, độ cao tăng tích lũy, số ô đường đi và
số cạnh đã thăm dò. Quy trình tạo 90.180 dòng dữ liệu chi tiết: ba bản ghi từ các
thuật toán không ngẫu nhiên và mười lăm bản ghi ngẫu nhiên trên mỗi bản đồ. Các bảng
tóm tắt trước tiên nhóm năm lần chạy ngẫu nhiên trong từng cặp bản đồ--thuật
toán rồi mới tổng hợp giữa các bản đồ, tránh xem năm seed trên cùng một địa hình
là năm mẫu địa hình độc lập. Thời gian được đo khi ba tiến trình chạy đồng thời
và được lưu cùng điều kiện thực thi này. Vì vậy, các giá trị thời gian chỉ phản
ánh lần chạy song song đã thực hiện và không được dùng để xếp hạng hoặc so sánh
tốc độ giữa các thuật toán. Mỗi thuật toán không ngẫu nhiên cũng chỉ có một phép
đo trên mỗi bản đồ. Muốn so sánh thời gian chạy, cần chạy tuần tự, lặp lại từng
cấu hình và báo cáo các thống kê ổn định như trung vị và khoảng tứ phân vị.

\subsection{Tái lập và chuẩn bị kiểm tra tính toàn vẹn}

Sinh dữ liệu, tạo ảnh, đánh giá thuật toán, phân tích và kiểm định được triển khai
thành các chương trình dòng lệnh riêng. Môi trường Conda và danh sách đầy đủ
phiên bản gói Python được cung cấp cùng mã nguồn. Công việc sinh dữ liệu và
đánh giá thuật toán dùng mã nhận dạng cấu hình cố định và có thể tiếp tục theo bộ ba hoặc theo
bản đồ hoàn chỉnh. Trình kiểm định cuối cùng kiểm tra cấu trúc cốt lõi, lưới
\texttt{grid} và \texttt{elevation}, khả năng tái tạo lưới va chạm và
\texttt{support\_elevation}, quy ước tọa độ, số lượng và liên thông nhiệm vụ,
điểm đầu--đích dùng chung và thứ tự độ khó. Tính hợp lệ của đường đi được kiểm
tra trong quá trình đánh giá. Ba danh sách SHA-256 được cung cấp cho toàn bộ
5.010 tệp bản đồ JSON, 5.010 tệp NPZ và siêu dữ liệu chính. Danh mục tệp hoàn chỉnh
ghi đường dẫn tương đối, loại nội dung và kích thước byte.

Hình~\ref{fig:terrain-triplets} minh họa cách cùng một mẫu địa hình và
cùng nhiệm vụ điểm đầu--đích được giữ qua ba mức Dễ--Trung bình--Khó trong mỗi bộ
ba ghép cặp.

\begin{figure}[htbp]
  \centering
  \IfFileExists{figures/terrain_triplet_001.png}{%
    \includegraphics[width=\linewidth,height=0.78\textheight,keepaspectratio]{figures/terrain_triplet_001.png}%
  }{%
    \figureplaceholder{75mm}{Tải lên
    \texttt{figures/terrain\_triplet\_001.png}.}%
  }
  \caption{\textbf{Các địa hình do hệ thống sinh ra ở ba mức độ khó ghép cặp.}
  Các cột lần lượt là Dễ, Trung bình và Khó; mỗi hàng là một họ địa hình.
  Các bản đồ trong một hàng dùng chung địa hình gốc và tọa độ điểm đầu--đích.
  Vùng xám biểu thị các vị trí mà tâm robot tham chiếu không thể đi qua. Hình
  được bố cục lại từ 15 ảnh do chương trình của dự án tạo ra; không thay đổi dữ
  liệu địa hình, vùng bị chặn hoặc tọa độ nhiệm vụ.}
  \label{fig:terrain-triplets}
\end{figure}

\section{Các bản ghi dữ liệu (Data Records)}

\subsection{Bản phát hành dữ liệu}

Bộ dữ liệu hoàn chỉnh được lưu trên \datasetrepository{} tại
\datasetdoi{}\cite{paired25ddata}. Mã nguồn tương ứng được lưu tại
\coderepository{}. Bản dữ liệu đầy đủ có dung
lượng khoảng 32,6 GB với 50.136 tệp trong danh mục. Kho mã nguồn dành cho phần
mềm và tài liệu dung lượng nhỏ, không thay thế kho lưu trữ lâu dài cho toàn bộ
dữ liệu.

\subsection{Tổ chức thư mục}

Thư mục gốc của bản phát hành là \texttt{dataset\_5010\_v1/}. Tệp bản đồ và ảnh
được chia theo họ địa hình để ba mức độ khó của một nhóm ghép cặp nằm gần nhau,
đồng thời tránh đặt hàng nghìn tệp vào một thư mục. Cấu trúc chính là:

\begin{verbatim}
dataset_5010_v1/
|-- maps/<family>/
|-- npz/<family>/
|-- images/
|   |-- terrain/<family>/
|   |-- analysis/<family>/
|   |-- paths/<family>/
|   `-- contact_sheets/
|-- metadata/
|   |-- dataset_summary.csv
|   |-- task_summary.csv
|   |-- splits/
|   |   |-- train.csv
|   |   |-- validation.csv
|   |   |-- test.csv
|   |   `-- split_protocol.json
|   |-- manifest.json
|   |-- validation_report.json
|   |-- pathfinding_benchmark_v1_trials.csv
|   |-- pathfinding_benchmark_v1_summary.csv
|   |-- pathfinding_benchmark_v1_protocol.json
|   |-- release_manifest.json
|   |-- release_filelist.csv
|   |-- npz_conversion.json
|   |-- npz_manifest.csv
|   |-- SHA256SUMS_CORE.txt
|   |-- SHA256SUMS_MAPS.txt
|   |-- SHA256SUMS_NPZ.txt
|   `-- benchmark_results/
|       |-- RESULTS.md
|       |-- benchmark_audit.json
|       |-- tables/
|       `-- figures/
|-- examples/load_map.py
|-- README.md
`-- SCHEMA.md
\end{verbatim}

Mỗi thư mục họ trong \texttt{maps/} có 1.002 tệp JSON: 334 mẫu gốc nhân với ba
độ khó. Tên tệp theo mẫu
\path{terrain_<family>_<index>_<difficulty>.json}; ví dụ
\path{terrain_mountain_001_hard.json}. Định danh dùng chung
\path{terrain_mountain_001} liên kết ba bản đồ Dễ, Trung bình và Khó.

\subsection{Nội dung và cấu trúc của tệp bản đồ JSON}

Mỗi tệp JSON chứa toàn bộ thông tin của một bản đồ, gồm định danh, seed, ma trận
độ cao thô, các lớp phân tích được tính từ độ cao, lưới va chạm cuối, mười nhiệm
vụ, quy ước tọa độ và điều hướng, thiết lập sinh dữ liệu, tiêu chí chấp nhận độ
khó và siêu dữ liệu của bản đồ. Mọi lớp raster có kích thước
$128\times128$. Mảng dùng chỉ số \texttt{[y][x]}, còn điểm đầu--đích của nhiệm vụ
dùng \texttt{[x,y]}. Trong lưới va chạm nhị phân, 0 là ô tâm robot hợp lệ và 1
là ô bị chặn.

Bảng~\ref{tab:raster-fields} tóm tắt các lớp raster chính. Độ cao thô được giữ
riêng với \texttt{support\_elevation}: lớp thứ nhất mô tả bề mặt tại ô, còn lớp
thứ hai là giá trị của mặt phẳng đỡ tại tâm vùng chiếm chỗ, được dùng để
tính độ dốc cạnh và chi phí đường đi.
Các mặt nạ nguyên nhân cho phép phân biệt vùng bị chặn do độ dốc, độ nhám,
chiều cao bậc hoặc chi phí đi qua địa hình tổng hợp.

\begin{table}[htbp]
  \centering
  \caption{Các lớp raster chính trong mỗi tệp JSON bản đồ.}
  \label{tab:raster-fields}
  \small
  \begin{tabular}{>{\ttfamily\raggedright\arraybackslash}p{0.28\linewidth}
                  >{\raggedright\arraybackslash}p{0.16\linewidth}
                  >{\raggedright\arraybackslash}p{0.48\linewidth}}
    \toprule
    \normalfont Trường & Mã hóa hoặc đơn vị & Ý nghĩa \\
    \midrule
    elevation & m & Độ cao bề mặt tại từng ô. \\
    grid & 0/1 & Lưới va chạm cuối đã xét vùng chiếm chỗ. \\
    point\_slope\_degrees & độ & Độ dốc gradient của bề mặt làm trơn nhẹ. \\
    support\_elevation & m & Giá trị của mặt phẳng đỡ bình phương tối thiểu tại tâm. \\
    footprint\_slope\_degrees & độ & Độ nghiêng của mặt phẳng đỡ trên vùng chiếm chỗ. \\
    roughness & m RMS & Biến thiên độ cao còn dư quanh mặt phẳng đỡ. \\
    step\_height & m & Chênh lệch độ cao kề nhau lớn nhất trong vùng chiếm chỗ. \\
    traversability\_cost & không thứ nguyên & Tổng có trọng số của độ dốc, độ nhám và chiều cao bậc đã chuẩn hóa. \\
    blocked\_by\_* & 0/1 & Mặt nạ nguyên nhân độ dốc, độ nhám, bậc và chi phí. \\
    centre\_blocked & 0/1 & Hợp của các mặt nạ nguyên nhân trước khi thêm biên an toàn. \\
    \bottomrule
  \end{tabular}
\end{table}

Mỗi nhiệm vụ gồm chỉ số, định danh riêng theo độ khó và định danh
ghép cặp, ô đầu và đích, nhóm khoảng cách, khoảng cách có trọng số tham chiếu,
độ khó tham chiếu và cờ nhiệm vụ chuẩn. Để tiện sử dụng, nhiệm vụ chuẩn còn
được lặp ở các trường cấp cao \texttt{start}, \texttt{goal},
\texttt{start\_z} và \texttt{goal\_z}. Từ điển trường đầy đủ, phép đổi tọa độ,
đơn vị và cách đọc tệp nằm trong \texttt{SCHEMA.md};
\texttt{examples/load\_map.py} cung cấp một ví dụ đọc tệp chỉ dùng thư viện
chuẩn, kèm các kiểm tra cấu trúc cơ bản.

\subsection{Các tệp raster NPZ}

Mỗi bản đồ JSON có một tệp cùng tên gốc trong \texttt{npz/<family>/}. Tệp này
chứa 13 trường raster ở Bảng~\ref{tab:raster-fields}, đều có kích thước
$128\times128$: lớp liên tục dùng \texttt{float32}, còn lưới và mặt nạ nhị phân
dùng \texttt{uint8}. Tệp không chứa đối tượng pickle và có thể đọc bằng
\texttt{numpy.load(path, allow\_pickle=False)}. NPZ là bản sao được tạo từ JSON
để truy cập mảng nhanh hơn; JSON vẫn là nguồn chuẩn cho định danh, nhiệm vụ,
thiết lập sinh và siêu dữ liệu. Quy trình chuyển đổi, đường dẫn, kích thước và mã băm từng tệp
nằm trong \texttt{npz\_conversion.json}, \texttt{npz\_manifest.csv} và
\texttt{SHA256SUMS\_NPZ.txt}.

\subsection{Siêu dữ liệu và kết quả đánh giá}

Tệp \texttt{metadata/dataset\_summary.csv} là chỉ mục bao quát cả năm họ, mỗi bản đồ
một dòng, chứa định danh, seed, giá trị thực tế của các tham số sinh, đường dẫn dữ liệu
và siêu dữ liệu của bản đồ. \texttt{metadata/task\_summary.csv} có một dòng cho mỗi
nhiệm vụ ghép cặp duy nhất, do đó có 16.700 dòng; định nghĩa nhiệm vụ chuẩn vẫn
nằm trong tệp JSON bản đồ. \texttt{metadata/manifest.json} ghi phạm vi phát hành, cấu
hình, bố cục ảnh và liên kết đến kết quả đánh giá.

Tệp kết quả chi tiết có 90.180 dòng. Các cột ghi định danh địa hình, nhóm ghép cặp,
nhiệm vụ, thuật toán, lần chạy, seed cố định, ngân sách, trạng thái thành công và hợp
lệ, thời gian chạy, chiều dài đường, chi phí có trọng số, độ cao tăng, số ô đường
đi và số cạnh đã thăm dò. CSV tóm tắt đi kèm nhóm lần chạy theo thuật toán, họ
và độ khó. Tệp JSON mô tả quy trình lưu quy tắc chuyển động, ngân sách thuật toán, số lần
chạy ngẫu nhiên,
phạm vi tạo ảnh và điều kiện thực thi, nhờ đó mỗi kết quả luôn đi kèm thiết lập
đã tạo ra nó.

Thư mục \texttt{benchmark\_results/} chứa bảng cấp bản đồ và bảng tổng hợp, bốn
hình tóm tắt kết quả, báo cáo cho người đọc và báo cáo kiểm tra có thể đọc
bằng máy. Các bản ghi
này phục vụ phần Kiểm định kỹ thuật và tái sử dụng sau này; chúng không thay thế
CSV lần chạy thô. Khả năng truy xuất nguồn gốc của bản phát hành được hỗ trợ bởi
\texttt{validation\_report.json}, ba danh sách SHA-256,
\texttt{release\_manifest.json} và \texttt{release\_filelist.csv}.

\subsection{Chỉ mục chia tập train, validation và test}

Ba chỉ mục cố định được cung cấp trong \texttt{metadata/splits/}. Việc chia tập
được thực hiện theo \texttt{matched\_group\_id}, không theo từng bản đồ, nên ba
bản đồ Dễ, Trung bình và Khó của cùng một địa hình gốc cùng mười nhiệm vụ
dùng chung luôn nằm trong một tập. Việc chia được thực hiện riêng trong từng họ địa hình.
Trong mỗi họ, 234 nhóm thuộc tập train, 50 nhóm thuộc validation và 50 nhóm
thuộc test; tổng cộng lần lượt là 1.170/250/250 nhóm, tương ứng
3.510/750/750 bản đồ.

Cách chia này có thể tái tạo. Trong từng họ, định danh nhóm được xếp theo giá trị
SHA-256 của tên quy trình, seed 20260830 và
\texttt{matched\_group\_id}; 234 nhóm đầu, 50 nhóm tiếp theo và 50 nhóm cuối lần
lượt được gán vào train, validation và test. Mỗi CSV có một dòng cho một nhóm,
kèm định danh và đường dẫn của cả ba bản đồ. Tệp
\texttt{split\_protocol.json} lưu quy tắc, seed, số lượng, mã băm của các tệp và
kết quả kiểm tra không rò rỉ nhóm. Các chỉ mục này dành cho phương pháp học;
phép đánh giá sáu thuật toán tham chiếu được chạy trên toàn bộ bản đồ và không huấn luyện
mô hình.

\subsection{Hình ảnh trong bản phát hành}

Với mỗi trong 5.010 bản đồ, \texttt{images/terrain/} chứa một ảnh địa hình và
\texttt{images/analysis/} chứa một ảnh phân tích khả năng đi qua địa hình, tổng cộng
10.020 ảnh bản đồ. Mỗi bản đồ còn có sáu ảnh kết quả, tương ứng với sáu thuật
toán lập kế hoạch, tạo thành 30.060 tệp dưới \texttt{images/paths/}. Với thuật
toán ngẫu nhiên, ảnh dùng lần chạy thành công đầu tiên theo thứ tự seed cố định;
nếu cả năm lần chạy đều không tìm được đường trong ngân sách, ảnh biểu diễn
trạng thái không có tuyến cuối cùng. Trạng thái tương ứng được lưu trong CSV
kết quả. Hai bảng ảnh tổng hợp cho phép xem nhanh các bộ ba địa hình và
kết quả trên bản đồ Khó. Các ảnh là hình minh họa được tạo từ dữ liệu để kiểm tra trực quan; tệp
JSON bản đồ và CSV lần chạy là nguồn dữ liệu số chính.

Hình~\ref{fig:analysis-products} cho thấy nội dung của ảnh phân tích, thay vì
chỉ liệt kê đường dẫn thư mục. Ba ví dụ đều là bản đồ mức Khó có chỉ số 001 và
được chọn trước theo họ, gồm \texttt{smooth\_obstacles}, \texttt{rugged} và
\texttt{plateau} để minh họa ba dạng địa hình khác nhau. Mỗi ảnh giữ cùng bố
cục ba ô hình của quy trình: trường độ cao và đường đồng mức; chi phí đi qua địa hình kèm vùng không thể đi;
và nguyên nhân từng tư thế tâm robot bị chặn.

\begin{figure}[htbp]
  \centering
  \textbf{(a) \texttt{smooth\_obstacles}}\par\smallskip
  \IfFileExists{figures/analysis_smooth_obstacles_hard.png}{%
    \includegraphics[width=0.98\linewidth]{figures/analysis_smooth_obstacles_hard.png}%
  }{\figureplaceholder{42mm}{Thiếu ảnh phân tích \texttt{smooth\_obstacles}.}}%
  \par\medskip
  \textbf{(b) Rugged}\par\smallskip
  \IfFileExists{figures/analysis_rugged_hard.png}{%
    \includegraphics[width=0.98\linewidth]{figures/analysis_rugged_hard.png}%
  }{\figureplaceholder{42mm}{Thiếu ảnh phân tích \texttt{rugged}.}}%
  \par\medskip
  \textbf{(c) Plateau}\par\smallskip
  \IfFileExists{figures/analysis_plateau_hard.png}{%
    \includegraphics[width=0.98\linewidth]{figures/analysis_plateau_hard.png}%
  }{\figureplaceholder{42mm}{Thiếu ảnh phân tích \texttt{plateau}.}}%
  \caption{\textbf{Kết quả phân tích địa hình trên ba họ minh họa ở mức Khó.}
  Các ô hình trong mỗi hàng lần lượt thể hiện độ cao, chi phí đi qua địa hình
  và nguyên nhân bị chặn sau khi xét vùng chiếm chỗ. Toàn bộ ảnh được tạo trực
  tiếp từ tệp JSON bản đồ có chỉ số 001; không chọn thủ công theo hình thức.}
  \label{fig:analysis-products}
\end{figure}

Hình~\ref{fig:planner-route-comparison} minh họa sáu loại ảnh kết quả được lưu
cho cùng một bản đồ và cùng tọa độ điểm đầu--đích. Hình này nhằm mô tả nội dung tệp,
không dùng để xếp hạng hoặc so sánh hiệu năng giữa các thuật toán.

\begin{figure}[htbp]
  \centering
  \begin{tabular}{cc}
    \includegraphics[width=0.46\textwidth]{figures/bfs_mountain_hard.png} &
    \includegraphics[width=0.46\textwidth]{figures/dijkstra_mountain_hard.png} \\
    \small (a) BFS & \small (b) Dijkstra \\[0.4em]
    \includegraphics[width=0.46\textwidth]{figures/a_star_mountain_hard.png} &
    \includegraphics[width=0.46\textwidth]{figures/prm_mountain_hard.png} \\
    \small (c) A* & \small (d) PRM \\[0.4em]
    \includegraphics[width=0.46\textwidth]{figures/rrt_connect_mountain_hard.png} &
    \includegraphics[width=0.46\textwidth]{figures/rrt_star_mountain_hard.png} \\
    \small (e) RRT-Connect & \small (f) RRT* \\
  \end{tabular}
  \caption{\textbf{Kết quả của sáu thuật toán trên cùng một nhiệm vụ.}
  Tất cả các ô hình dùng \texttt{terrain\_mountain\_001\_hard} và cùng tọa độ
  điểm đầu--đích. Các tuyến được hiển thị là kết quả do chương trình đánh giá tạo ra và
  không đi qua ô bị chặn.}
  \label{fig:planner-route-comparison}
\end{figure}

\section{Kiểm định kỹ thuật (Technical Validation)}

\subsection{Kiểm định cấu trúc cốt lõi và thiết kế ghép cặp}

Trình kiểm định phát hành đọc lại toàn bộ 5.010 tệp JSON bản đồ thay vì chỉ kiểm
tra siêu dữ liệu tóm tắt. Các kiểm tra bao gồm cấu trúc cốt lõi, hình dạng và miền
giá trị của \texttt{grid} và \texttt{elevation}, khả năng tái tạo lưới va chạm
và \texttt{support\_elevation}, quy ước tọa độ \texttt{[y][x]} và
\texttt{[x,y]}, số lượng nhiệm vụ, trạng thái hợp lệ của điểm đầu--đích, khả năng
liên thông theo mô hình chuyển động công bố và sự nhất quán của các tọa độ nhiệm
vụ trong mỗi bộ ba. Báo cáo
\texttt{validation\_report.json} ghi nhận 5.010/5.010 bản đồ và 50.100/50.100 bản
ghi bản đồ--nhiệm vụ hợp lệ, không có lỗi hoặc cảnh báo.

Toàn bộ 5.010 tệp NPZ còn được kiểm tra CRC của ZIP, đúng tập 13 trường, phần
tiêu đề NPY, kích thước $128\times128$, kiểu dữ liệu và độ dài phần dữ liệu. Báo cáo chuyển
đổi không ghi nhận tệp lỗi.

Một kiểm tra bổ sung trên siêu dữ liệu và kết quả đánh giá xem xét 1.670
nhóm ghép cặp. Cả 1.670 nhóm đều có hệ số cường độ địa hình $q$ tăng nghiêm ngặt và tỷ lệ ô
đi được giảm nghiêm ngặt theo thứ tự Dễ--Trung bình--Khó; số vi phạm của
mỗi điều kiện bằng 0. Kiểm tra này quan trọng vì ba bản đồ trong một bộ ba là
mẫu ghép cặp và không được xem như ba mẫu độc lập.

\subsection{Kiểm định phân tầng độ khó}

Hình~\ref{fig:difficulty-overview} mở rộng kiểm tra thứ tự bằng các thống kê mô
tả trên toàn bộ 5.010 bản đồ. Mỗi điểm là trung bình của 334 bản đồ trong một tổ hợp
họ--độ khó và thanh sai số là một độ lệch chuẩn mẫu giữa các bản đồ. Từ Dễ đến
Khó, cả năm họ đều giảm tỷ lệ ô đi được đồng thời tăng hệ số $q$ và
độ dốc trung bình trên vùng chiếm chỗ. Hình được dùng để xác nhận rằng các nhãn độ
khó tuân theo các tiêu chí hiệu chỉnh đã công bố, không dùng để so sánh hiệu năng
thuật toán.

\begin{figure}[p]
  \centering
  \IfFileExists{figures/dataset_difficulty_overview.png}{%
    \includegraphics[width=\linewidth,height=0.68\textheight,keepaspectratio]{figures/dataset_difficulty_overview.png}%
  }{%
    \figureplaceholder{90mm}{Tải lên
    \texttt{figures/dataset\_difficulty\_overview.png}.}%
  }
  \caption{\textbf{Hiệu chỉnh độ khó địa hình ghép cặp trên toàn bộ bộ dữ liệu.}
  Bốn ô hình lần lượt trình bày tỷ lệ ô đi được, hệ số cường độ địa
  hình $q$, độ dốc trung bình trên vùng chiếm chỗ và số thành phần liên thông theo năm họ và ba mức
  độ khó. Điểm và thanh sai số là trung bình $\pm$ độ lệch chuẩn mẫu trên 334
  bản đồ cho mỗi tổ hợp họ--độ khó.}
  \label{fig:difficulty-overview}
\end{figure}

\subsection{Kiểm định độ khó lập kế hoạch theo cặp}

Để kiểm tra độ khó bằng đại lượng không tham gia vào quy tắc chấp nhận nhãn,
chúng tôi dùng nhiệm vụ dài chuẩn có cùng tọa độ điểm đầu--đích trong ba bản đồ
của mỗi bộ ba. Đơn vị thống kê là một bộ ba ghép cặp; mỗi họ có $n=334$ bộ ba.
Với các thuật toán ngẫu nhiên, năm lần chạy trước hết được gộp thành tỷ lệ thành
công cấp bản đồ. Sau đó, với từng chỉ số, chênh lệch
$\Delta=\mathrm{Khó}-\mathrm{Dễ}$ được tính trong mỗi bộ ba.
Bảng~\ref{tab:paired-planning-effects} báo cáo
trung bình ghép cặp và khoảng tin cậy 95\% xấp xỉ chuẩn
$\bar\Delta\pm1{,}96s_\Delta/\sqrt{n}$.

\begin{table}[htbp]
  \centering
  \caption{Thay đổi độ khó lập kế hoạch từ Dễ sang Khó. $\Delta C^*$ là thay đổi
  chi phí đường đi có trọng số tối ưu của A*. Ba cột còn lại là thay đổi tỷ lệ
  thành công, tính theo điểm phần trăm. Mỗi ô là trung bình ghép cặp [95\% CI]
  trên 334 bộ ba.}
  \label{tab:paired-planning-effects}
  \footnotesize
  \setlength{\tabcolsep}{4pt}
  \renewcommand{\arraystretch}{1.15}
  \begin{tabular}{lcccc}
    \toprule
    Họ địa hình & $\Delta C^*$ & $\Delta$ PRM & $\Delta$ RRT-Connect & $\Delta$ RRT* \\
    \midrule
    \texttt{smooth\_obstacles} & \shortstack{7,45\\{[7,02; 7,87]}} &
      \shortstack{$-0{,}06$\\{$[-0{,}18; 0{,}06]$}} &
      \shortstack{0,00\\{[0,00; 0,00]}} & \shortstack{0,00\\{[0,00; 0,00]}} \\
    \texttt{rolling} & \shortstack{26,65\\{[25,05; 28,25]}} &
      \shortstack{$-39{,}64$\\{$[-43{,}78; -35{,}50]$}} &
      \shortstack{$-22{,}22$\\{$[-25{,}85; -18{,}58]$}} &
      \shortstack{$-27{,}90$\\{$[-31{,}80; -24{,}01]$}} \\
    \texttt{mountain} & \shortstack{26,04\\{[24,43; 27,65]}} &
      \shortstack{$-44{,}85$\\{$[-48{,}90; -40{,}80]$}} &
      \shortstack{$-25{,}57$\\{$[-29{,}16; -21{,}98]$}} &
      \shortstack{$-29{,}52$\\{$[-33{,}33; -25{,}71]$}} \\
    \texttt{rugged} & \shortstack{33,62\\{[31,87; 35,36]}} &
      \shortstack{$-54{,}31$\\{$[-58{,}45; -50{,}17]$}} &
      \shortstack{$-34{,}67$\\{$[-38{,}66; -30{,}69]$}} &
      \shortstack{$-42{,}46$\\{$[-46{,}75; -38{,}16]$}} \\
    \texttt{plateau} & \shortstack{4,37\\{[3,98; 4,75]}} &
      \shortstack{0,00\\{[0,00; 0,00]}} &
      \shortstack{0,00\\{[0,00; 0,00]}} & \shortstack{0,00\\{[0,00; 0,00]}} \\
    \bottomrule
  \end{tabular}
\end{table}

Chi phí tối ưu tăng từ Dễ sang Khó ở cả năm họ, với khoảng tin cậy không chứa
0. Tỷ lệ thành công dưới ngân sách cố định giảm rõ ở \texttt{rolling},
\texttt{mountain} và \texttt{rugged}, nhưng bão hòa gần 100\% ở
\texttt{smooth\_obstacles} và \texttt{plateau}. Vì vậy, phân tầng độ khó làm
tăng chi phí lập kế hoạch một cách nhất quán, còn tác động lên xác suất thành
công phụ thuộc họ địa hình và thuật toán; nhãn không được diễn giải như một bảo
đảm rằng mọi thuật toán đều suy giảm trên mọi họ.

\subsection{Kiểm định đường đi và kết quả đánh giá}

Mỗi đường đi do thuật toán trả về được kiểm tra trước khi lưu kết quả lần chạy.
Phép kiểm tra xét ô đầu--đích, va chạm, quy tắc cắt góc chéo, giới hạn độ dốc
cạnh và tính liên tục tám hướng sau khi chuyển đường thành chuỗi ô. Trong 90.180 lần chạy, 83.998 lần
trả về đường hợp lệ và 6.182 lần không tìm được đường trong ngân sách cố định;
không có đường nào được đánh dấu thành công nhưng không hợp lệ. Cả 15.030 lần
chạy của BFS, Dijkstra và A* đều thành công trên 5.010 nhiệm vụ chuẩn, xác nhận
rằng mỗi nhiệm vụ chuẩn có ít nhất một đường đi trên lưới tham chiếu; do đó
không có nhiệm vụ nào mà cả sáu thuật toán đều thất bại. Các lần không tìm được
đường chỉ xuất hiện ở PRM, RRT-Connect và RRT* với lấy mẫu
ngẫu nhiên và ngân sách hữu hạn; chúng không có nghĩa là nhiệm vụ không liên
thông. Các số này chỉ kiểm tra tính hợp lệ và xác nhận rằng phép đánh giá ghi nhận
được cả lần chạy thành công lẫn lần không tìm được đường dưới ngân sách cố định;
chúng không được dùng để xếp hạng thuật toán.

Kiểm tra phát hành đọc các cờ \texttt{success} và \texttt{path\_valid} trong
90.180 dòng lần chạy, đồng thời kiểm kê 30.060 cặp bản đồ--thuật toán. Kết quả
không có dòng \texttt{success=True} nào mang cờ \texttt{path\_valid} khác
\texttt{True}, và không thiếu ảnh kết quả nào. Ngoài ra, A* và Dijkstra cho cùng
chi phí có trọng số đến sáu chữ số
thập phân được lưu trên cả 5.010 nhiệm vụ chuẩn; chênh lệch tuyệt đối lớn nhất
trong tệp kết quả là 0. Đây là kiểm tra chéo giữa hai cách triển khai tìm kiếm tối
ưu trên cùng mô hình chi phí.

Bảng~\ref{tab:validation-summary} tóm tắt phạm vi và kết quả của các kiểm tra
kỹ thuật chính.

\begin{table}[htbp]
  \centering
  \caption{Tóm tắt các kiểm tra kỹ thuật của bản phát hành.}
  \label{tab:validation-summary}
  \small
  \begin{tabular}{L{0.43\linewidth}r r}
    \toprule
    Nội dung xác minh & Số đối tượng kiểm tra & Số lỗi phát hiện \\
    \midrule
    Cấu trúc cốt lõi của tệp bản đồ & 5.010 bản đồ & 0 \\
    Tính hợp lệ của nhiệm vụ điểm đầu--đích & 50.100 bản ghi bản đồ--nhiệm vụ & 0 \\
    Thứ tự hệ số $q$ Dễ--Trung bình--Khó & 1.670 bộ ba & 0 \\
    Thứ tự tỷ lệ ô đi được & 1.670 bộ ba & 0 \\
    Rò rỉ \texttt{matched\_group\_id} giữa các tập & 1.670 bộ ba & 0 \\
    Tính hợp lệ của đường được báo thành công & 83.998 đường & 0 \\
    Tính đầy đủ của ảnh kết quả thuật toán & 30.060 ảnh kỳ vọng & 0 \\
    Sự nhất quán chi phí A*--Dijkstra & 5.010 nhiệm vụ chuẩn & 0 \\
    \bottomrule
  \end{tabular}
\end{table}

Ở hàng kiểm tra ảnh kết quả, giá trị 0 trong cột ``Số lỗi phát hiện'' nghĩa là cả
30.060 ảnh kỳ vọng đều tồn tại. Sự tồn tại của ảnh không đồng nghĩa thuật toán
luôn tìm được đường; trạng thái thành công hoặc không tìm được đường phải được
đọc từ CSV lần chạy.
Kiểm tra cách chia tập còn xác nhận mỗi trong 1.670 nhóm xuất hiện đúng một lần
trên ba chỉ mục và mỗi tập có đúng số nhóm quy định cho từng họ địa hình.

\subsection{Khả năng tái lập}

Các seed và mã nhận dạng cấu hình được giữ trong bản phát hành dữ liệu; phiên bản
Python và thư viện được cung cấp trong kho mã nguồn. Bộ kiểm thử hồi quy kiểm tra
mặt phẳng ngang, khôi phục độ dốc của mặt
phẳng đỡ, khả năng tái tạo và cách chia nhiệm vụ theo khoảng cách, phép chuyển cạnh thành chuỗi ô đối xứng, quy tắc
cắt góc, sử dụng trọng số chi phí và sự đồng thuận A*/Dijkstra. Ba danh sách
SHA-256 cho phép phát hiện thay đổi ở 5.010 bản đồ JSON, các tệp NPZ và siêu dữ liệu
chính sau khi tải.

\section{Lưu ý sử dụng (Usage Notes)}

Các lớp raster dùng chỉ số \texttt{[y][x]}, còn ô đầu--đích được lưu theo
\texttt{[x,y]}; đảo hai quy ước này sẽ làm đổi hàng thành cột mà không nhất thiết
gây lỗi đọc tệp. Tệp \texttt{examples/load\_map.py} minh họa cách đọc JSON đúng
chỉ bằng thư viện chuẩn. Với xử lý số,
tệp NPZ cùng tên gốc cho phép tải trực tiếp các mảng raster; siêu dữ liệu nhiệm vụ và
cấu hình vẫn cần đọc từ JSON. Nhiệm vụ
\texttt{long} được chọn trên bản đồ Khó từ dải 0,75--0,98 của các khoảng
cách có trọng số đã xếp hạng rồi sao chép sang hai bản đồ Dễ và Trung bình ghép
cặp. Nhiệm vụ dài nhất được chọn làm nhiệm vụ chuẩn.

Lưới va chạm và các lớp về khả năng đi qua địa hình được phát hành cho vùng chiếm chỗ hình tròn
bán kính 0,45 m cùng các giới hạn hình học đã công bố. Khi đánh giá robot khác,
người dùng cần tính lại các lớp có xét vùng chiếm chỗ từ \texttt{elevation}, cập nhật
dải biên an toàn và quy tắc chuyển động, rồi kiểm tra lại tính liên thông của
nhiệm vụ. Làm giãn trực tiếp \texttt{grid} đã phát hành sẽ tính vùng chiếm chỗ tham
chiếu hai lần, trừ khi chủ đích là thêm một biên an toàn riêng.

Với phương pháp học, cần tuân theo các chỉ mục chia tập cố định và xem
\texttt{matched\_group\_id} là đơn vị không thể tách. Trong phân tích ghép cặp,
một bộ ba là đơn vị thống kê; ba mức độ khó và các seed lặp của thuật toán ngẫu
nhiên không phải các mẫu địa hình độc lập. Bộ dữ liệu chỉ mô tả hình học tĩnh
dưới dạng trường độ cao 2,5D. Nó không mô hình hóa biến dạng đất, trượt bánh và
động lực học xe; thực vật, nước, bùn hoặc tính chất ngữ nghĩa của bề mặt; nhiễu
cảm biến và bất định định vị; cầu, phần nhô ra hay nhiều tầng bề mặt tại cùng
một tọa độ. Mô hình chuyển động tám hướng cũng không biểu diễn hướng thân xe,
bán kính quay, cơ cấu lái hoặc động học riêng của từng robot. Toàn bộ dữ liệu
được tạo ở độ phân giải 0,25 m mỗi ô; ảnh hưởng của việc đổi độ phân giải chưa
được kiểm định. Do địa hình là dữ liệu tổng hợp, mức độ tương đồng và khả năng
chuyển giao sang địa hình thực cũng chưa được xác nhận định lượng. Vì vậy, các
ngưỡng robot và kết quả của các thuật toán tham chiếu không nên được diễn giải như giới hạn
vận hành hoặc bảng xếp hạng phổ quát.

\section{Tính sẵn có của dữ liệu (Data Availability)}

Bản phát hành đầy đủ gồm tệp JSON bản đồ chuẩn, raster NPZ không dùng pickle,
nhiệm vụ điểm đầu--đích, chỉ mục chia tập cố định, ảnh được tạo từ dữ liệu, kết quả đánh giá,
siêu dữ liệu, báo cáo kiểm định, danh mục tệp và mã kiểm tra SHA-256 được lưu tại
\datasetrepository{} và định danh bằng
\datasetdoi{}\cite{paired25ddata}. Bộ dữ liệu được phát hành theo
\datasetlicense{}. Kho mã nguồn không thay thế kho dữ liệu lưu trữ
lâu dài cho toàn bộ 32,6 GB dữ liệu.

\section{Tính sẵn có của mã nguồn (Code Availability)}

Mã Python dùng để sinh địa hình, phân tích khả năng đi qua địa hình, chọn nhiệm vụ,
chạy sáu thuật toán, tạo ảnh, kiểm định và chuẩn bị bản phát hành được cung
cấp tại \coderepository{}. Kho mã nguồn chứa phiên bản tương ứng với bản thảo
đánh giá này, môi trường Conda, \texttt{requirements-lock.txt}, kiểm thử hồi
quy và hướng dẫn tái lập. Mã nguồn được cấp phép theo giấy phép MIT.

\begin{thebibliography}{99}

\bibitem{rugd}
Wigness, M. B., Eum, S., Rogers, J. G. III, Han, D. K. \& Kwon, H. A RUGD
dataset for autonomous navigation and visual perception in unstructured outdoor
environments. In \textit{2019 IEEE/RSJ International Conference on Intelligent
Robots and Systems (IROS)}, 5000--5007 (2019).
\url{https://doi.org/10.1109/IROS40897.2019.8968283}.

\bibitem{rellis3d}
Jiang, P., Osteen, P. R., Wigness, M. B. \& Saripalli, S. RELLIS-3D dataset:
data, benchmarks and analysis. In \textit{2021 IEEE International Conference on
Robotics and Automation (ICRA)}, 1110--1116 (2021).
\url{https://doi.org/10.1109/ICRA48506.2021.9561251}.

\bibitem{wildscenes}
Vidanapathirana, K. et al. WildScenes: a benchmark for 2D and 3D semantic
segmentation in large-scale natural environments. \textit{The International
Journal of Robotics Research} \textbf{44}, 532--549 (2025).
\url{https://doi.org/10.1177/02783649241278369}.

\bibitem{goose}
Mortimer, P. et al. The GOOSE dataset for perception in unstructured
environments. In \textit{2024 IEEE International Conference on Robotics and
Automation (ICRA)}, 14838--14844 (2024).
\url{https://doi.org/10.1109/ICRA57147.2024.10611298}.
Data and code: \url{https://github.com/FraunhoferIOSB/goose_dataset}.

\bibitem{orfd}
Min, C. et al. ORFD: a dataset and benchmark for off-road freespace detection.
In \textit{2022 International Conference on Robotics and Automation (ICRA)},
2532--2538 (2022).
\url{https://doi.org/10.1109/ICRA46639.2022.9812139}.

\bibitem{tartandrive}
Triest, S., Sivaprakasam, M., Wang, S. J., Wang, W., Johnson, A. M. \& Scherer,
S. TartanDrive: a large-scale dataset for learning off-road dynamics models. In
\textit{2022 International Conference on Robotics and Automation (ICRA)}
(2022). \url{https://doi.org/10.1109/ICRA46639.2022.9811648}.

\bibitem{stone}
Park, K. et al. STONE dataset: a scalable multi-modal surround-view 3D
traversability dataset for off-road robot navigation. \textit{arXiv preprint}
arXiv:2603.09175 (2026).
\url{https://doi.org/10.48550/arXiv.2603.09175}.

\bibitem{vertibench}
Xu, T. et al. Verti-Bench: A general and scalable off-road mobility benchmark
for vertically challenging terrain. In \textit{Robotics: Science and Systems}
(2025). \url{https://doi.org/10.15607/RSS.2025.XXI.138}.
Project, data and code: \url{https://people.cs.gmu.edu/~xiao/Research/Verti-Bench/}.

\bibitem{planetarypathbench}
Chanc\'{a}n, M., Banerjee, A. \& Nikolakopoulos, G. Planetary terrain datasets
and benchmarks for rover path planning. \textit{arXiv preprint}
arXiv:2512.21438 (2025).
\url{https://doi.org/10.48550/arXiv.2512.21438}.
Code and data: \url{https://github.com/mchancan/PlanetaryPathBench}.

\bibitem{barn}
Perille, D., Truong, A., Xiao, X. \& Stone, P. Benchmarking metric ground
navigation. In \textit{2020 IEEE International Symposium on Safety, Security,
and Rescue Robotics (SSRR)}, 116--121 (2020).
\url{https://doi.org/10.1109/SSRR50563.2020.9292572}.

\bibitem{fair}
Wilkinson, M. D. et al. The FAIR Guiding Principles for scientific data
management and stewardship. \textit{Sci. Data} \textbf{3}, 160018 (2016).
\url{https://doi.org/10.1038/sdata.2016.18}.

\bibitem{numpy}
Harris, C. R. et al. Array programming with NumPy. \textit{Nature}
\textbf{585}, 357--362 (2020).
\url{https://doi.org/10.1038/s41586-020-2649-2}.

\bibitem{scipy}
Virtanen, P. et al. SciPy 1.0: fundamental algorithms for scientific computing
in Python. \textit{Nature Methods} \textbf{17}, 261--272 (2020).
\url{https://doi.org/10.1038/s41592-019-0686-2}.

\bibitem{dijkstra}
Dijkstra, E. W. A note on two problems in connexion with graphs.
\textit{Numerische Mathematik} \textbf{1}, 269--271 (1959).
\url{https://doi.org/10.1007/BF01386390}.

\bibitem{astar}
Hart, P. E., Nilsson, N. J. \& Raphael, B. A formal basis for the heuristic
determination of minimum cost paths. \textit{IEEE Transactions on Systems
Science and Cybernetics} \textbf{4}, 100--107 (1968).
\url{https://doi.org/10.1109/TSSC.1968.300136}.

\bibitem{prm}
Kavraki, L. E., Svestka, P., Latombe, J.-C. \& Overmars, M. H. Probabilistic
roadmaps for path planning in high-dimensional configuration spaces.
\textit{IEEE Transactions on Robotics and Automation} \textbf{12}, 566--580
(1996). \url{https://doi.org/10.1109/70.508439}.

\bibitem{rrtconnect}
Kuffner, J. J. \& LaValle, S. M. RRT-Connect: An efficient approach to
single-query path planning. In \textit{Proceedings of the IEEE International
Conference on Robotics and Automation}, 995--1001 (2000).
\url{https://doi.org/10.1109/ROBOT.2000.844730}.

\bibitem{rrtstar}
Karaman, S. \& Frazzoli, E. Sampling-based algorithms for optimal motion
planning. \textit{The International Journal of Robotics Research} \textbf{30},
846--894 (2011). \url{https://doi.org/10.1177/0278364911406761}.

\bibitem{paired25ddata}
Tran, Q. M. A paired multi-difficulty 2.5D terrain dataset for ground-robot
path planning. \textit{Zenodo} (2026).
\url{https://doi.org/10.5281/zenodo.22074838}.

\end{thebibliography}

\section*{Đóng góp của tác giả (Author Contributions)}

T.Q.M. đề xuất ý tưởng và thiết kế bộ dữ liệu; phát triển phần mềm sinh địa
hình, phân tích khả năng đi qua địa hình, lựa chọn nhiệm vụ, đánh giá các thuật toán
và kiểm định; sinh, kiểm định, quản lý và lập tài liệu cho bản phát hành dữ liệu;
thực hiện phân tích, chuẩn bị hình và bảng; soạn thảo, chỉnh sửa và phê duyệt
bản thảo.

\section*{Xung đột lợi ích (Competing Interests)}

Tác giả tuyên bố không có xung đột lợi ích.

\section*{Tài trợ (Funding)}

Nghiên cứu này không nhận tài trợ từ nguồn bên ngoài.

\end{document}
